\documentclass[12pt]{article}
\usepackage[utf8]{inputenc}\usepackage[T1]{fontenc}
\usepackage[a4paper,margin=2.5cm]{geometry}
\usepackage{amsmath,amssymb,amsthm,mathtools,mathrsfs}
\usepackage[dvipsnames]{xcolor}
\usepackage{booktabs,enumitem,hyperref}
\setlist{nosep,leftmargin=1.8em}
\hypersetup{colorlinks=true,linkcolor=blue!60!black,citecolor=green!40!black,urlcolor=blue!50!black}

\theoremstyle{plain}
\newtheorem{theorem}{Theorem}[section]
\newtheorem{proposition}[theorem]{Proposition}
\newtheorem{corollary}[theorem]{Corollary}
\newtheorem{lemma}[theorem]{Lemma}
\newtheorem*{observation}{Unifying Observation}
\theoremstyle{definition}
\newtheorem{definition}[theorem]{Definition}
\newtheorem{axiom}[theorem]{Axiom}
\newtheorem*{conjecture}{Conjecture}
\theoremstyle{remark}
\newtheorem{remark}[theorem]{Remark}

\newcommand{\Poinc}{\mathrm{Poinc}}
\newcommand{\SO}{\mathrm{SO}}
\newcommand{\SU}{\mathrm{SU}}
\newcommand{\SL}{\mathrm{SL}}
\newcommand{\ISO}{\mathrm{ISO}}
\newcommand{\Diff}{\mathrm{Diff}}
\newcommand{\Ker}{\mathrm{Ker}}
\newcommand{\Isom}{\mathrm{Isom}}
\newcommand{\dd}{\mathrm{d}}

\title{\textbf{The Equations of Physics from Orbital Structure:\\Quantum Field Theory, the Unifying Observation,\\and Constraints on Quantum Gravity}}

\author{Jocelyn Nemb\'e\thanks{Email: jnembe777@gmail.com, jnembe@root-insights.org}\\[0.3em]
\small Roots Insights Laboratory, Singapore\\
\small National Institute of Management Sciences, Libreville, Gabon}

\date{}

\begin{document}
\maketitle

% ============================================================================
\begin{abstract}
This paper completes the ``Physics from Orbital Structure'' series by (I)~deriving the structure of quantum field theory --- the existence of quantum fields, the spin-statistics connection, the CPT theorem, and gauge invariance --- from the same three-step orbital pattern used for special relativity (Paper~1), general relativity (Paper~2), and quantum mechanics (Paper~3); and (II)~formalizing the observation that all four theories follow the same pattern as a \textbf{Unifying Observation} (not a theorem in the strict sense, but an empirical pattern across four rigorously proved instances).

For QFT, the symmetry group is the Poincar\'e group acting on Fock space, the conservation principle is unitarity of the S-matrix combined with the \textbf{cluster decomposition principle}, and the uniqueness is provided by representation theory (following Weinberg's program). We show that (i)~the cluster decomposition principle is the multi-particle analog of the Bianchi identity (GR) and the Casimir eigenvalue (QM); (ii)~massless spin-2 particles force general covariance, closing a self-consistency loop with Paper~2; (iii)~a Lagrangian exists \emph{a posteriori} in every case (Cartan--Vermeil--Weyl for GR, Helmholtz for QM, Weinberg reconstruction for QFT); (iv)~seven constraints on quantum gravity follow from the empty corner of the classification grid $(\kappa, \rho, \sigma)$.

We propose the \textbf{Orbital Integrability Conjecture}: the equations of physics are integrability conditions for the self-consistency of orbital structure.
\end{abstract}

\noindent\textbf{Keywords:} unification, orbital structure, quantum field theory, spin-statistics, cluster decomposition, Bianchi--Casimir--Cluster correspondence, quantum gravity constraints, Lagrangian-free foundations

\noindent\textbf{MSC 2020:} 81T99, 83E99, 22E70, 53Z05


% ============================================================================
\section{Introduction}
\label{sec:intro}
% ============================================================================

\subsection{Summary of Papers 1--3}

Papers~1--3 of this series~\cite{nembe2026p1,nembe2026p2,nembe2026p3a} derived the fundamental equations of special relativity, general relativity, and quantum mechanics from a single three-step pattern:
\begin{center}
\fbox{\textbf{Symmetry group} $\to$ \textbf{Conservation principle} $\to$ \textbf{Uniqueness theorem} $\to$ \textbf{The equation}}
\end{center}
No Lagrangian was used as input in any of the three derivations. The theories are classified by a pair $(\kappa, \rho)$: SR $= (10, \alpha)$; GR $= (<\!10, \alpha)$; QM $= (10, \beta)$.

\subsection{What this paper adds}

This capstone paper has two parts. \textbf{Part~I} (\S\S\ref{sec:fock}--\ref{sec:gauge}) extends the orbital pattern to quantum field theory, the multi-particle extension of QM. Following Weinberg~\cite{weinberg1995}, quantum fields \emph{emerge} from particles (Poincar\'e representations), unitarity, and the cluster decomposition principle --- no classical field is quantized. \textbf{Part~II} (\S\S\ref{sec:framework}--\ref{sec:oip}) formalizes the unifying observation, proves the structural correspondences, establishes the self-consistency loop, shows the Lagrangian is always a consequence, and derives constraints on quantum gravity.

We note that the QFT derivation (Part~I) is essentially a reformulation of Weinberg's program~\cite{weinberg1995} in the orbital language. The genuinely new contributions are: the Bianchi--Casimir--Cluster correspondence (\S\ref{sec:correspondences}), the self-consistency loop (\S\ref{sec:loop}), and the QG constraints (\S\ref{sec:qg}).


% ============================================================================
% PART I: QFT FROM ORBITAL STRUCTURE
% ============================================================================

\part*{Part I: Quantum Field Theory from Orbital Structure}

\section{From Particles to Fock Space}
\label{sec:fock}
% ============================================================================

\begin{definition}[Multi-particle states]\label{def:fock}
Given a single-particle Hilbert space $\mathcal{H}_1$ (an irreducible unitary representation of the Poincar\'e group with mass $m$ and spin $s$, cf.\ Paper~3a), the \textbf{Fock space} is:
\begin{equation}
\mathcal{F} = \bigoplus_{n=0}^{\infty} \mathcal{H}_n, \qquad \mathcal{H}_0 = \mathbb{C}\,|\Omega\rangle \;\text{(vacuum)},
\end{equation}
where $\mathcal{H}_n$ is the $n$-particle Hilbert space (to be specified as symmetric or antisymmetric tensor product --- this will be \emph{determined} by the theory, not postulated).
\end{definition}

\begin{definition}[Creation and annihilation]\label{def:creation}
For each momentum $\mathbf{p}$ and spin component $\sigma$: $a^\dagger(\mathbf{p}, \sigma)$ creates a particle; $a(\mathbf{p}, \sigma)$ destroys one. The Poincar\'e group acts by:
\begin{equation}\label{eq:poincare-on-a}
U(\Lambda, a)\,a^\dagger(\mathbf{p}, \sigma)\,U^{-1}(\Lambda, a) = e^{i(\Lambda p)\cdot a}\sqrt{\frac{(\Lambda p)^0}{p^0}} \sum_{\bar\sigma} D^{(s)}_{\bar\sigma\sigma}(W(\Lambda, p))\,a^\dagger(\Lambda\mathbf{p}, \bar\sigma),
\end{equation}
where $W(\Lambda, p) = L^{-1}_{\Lambda p}\Lambda L_p$ is the Wigner rotation (an element of the little group) and $L_p$ is a standard boost~\cite{weinberg1995}.
\end{definition}

\begin{axiom}[Poincar\'e invariance of the S-matrix]\label{ax:QF1}
The scattering operator $S$ commutes with all Poincar\'e transformations: $[U(\Lambda, a), S] = 0$.
\end{axiom}

\begin{axiom}[Unitarity]\label{ax:QF2}
The S-matrix preserves probability: $S^\dagger S = SS^\dagger = 1$.
\end{axiom}


% ============================================================================
\section{Cluster Decomposition and the Emergence of Fields}
\label{sec:cluster}
% ============================================================================

\begin{axiom}[Cluster decomposition]\label{ax:QF3}
If two sets of particles $A$ and $B$ are prepared at spacelike separation $d$ with $d \to \infty$:
\begin{equation}\label{eq:cluster}
S_{A \cup B} \to S_A \otimes S_B.
\end{equation}
The scattering of widely separated subsystems factorizes: distant experiments are statistically independent.
\end{axiom}

\begin{remark}[Physical content]
Cluster decomposition is the mathematical expression of \emph{locality}: what happens in one laboratory does not instantaneously affect a distant laboratory. It is the multi-particle analog of the single-particle conservation principle (Paper~3a): both ensure that the physical structure is ``locally self-consistent.''
\end{remark}

\begin{theorem}[Quantum fields from cluster decomposition --- Weinberg~\cite{weinberg1995}]\label{thm:fields-from-cluster}
Axioms QF1--QF3 imply that the interaction Hamiltonian can be written as:
\begin{equation}\label{eq:H-local}
V = \int \dd^3x\,\mathcal{H}(\mathbf{x}, t),
\end{equation}
where $\mathcal{H}(x)$ is a \emph{local} density constructed from \textbf{quantum fields}:
\begin{equation}\label{eq:field}
\phi_l(x) = \sum_\sigma \int \frac{\dd^3p}{(2\pi)^{3/2}\sqrt{2p^0}} \left[u_l(\mathbf{p}, \sigma)\,e^{ip\cdot x}\,a(\mathbf{p}, \sigma) + v_l(\mathbf{p}, \sigma)\,e^{-ip\cdot x}\,a^\dagger(\mathbf{p}, \sigma)\right],
\end{equation}
where $u_l$ and $v_l$ are coefficient functions determined by the Lorentz transformation properties of the field. The field $\phi_l(x)$ transforms under a finite-dimensional representation of the Lorentz group.
\end{theorem}

\begin{proof}[Proof outline]
The proof follows Weinberg~\cite{weinberg1995}, Chapters~3--4. Write $S = \mathcal{T}\exp(-i\int V(t)\,\dd t)$ with $V(t) = \int \dd^3x\,\mathcal{H}(x)$. The connected part of the S-matrix, $S^c$, satisfies cluster decomposition iff $\mathcal{H}(x)$ is built from products of \emph{local} operators. Poincar\'e covariance of $S$ requires $\mathcal{H}(x)$ to transform as a Lorentz scalar under $U(\Lambda, a)$. The most general Lorentz-scalar local operator built from creation and annihilation operators is a polynomial in fields of the form~\eqref{eq:field}. The coefficient functions $u_l, v_l$ are determined by requiring $\phi_l(x)$ to transform under a definite finite-dimensional representation $(j_1, j_2)$ of $\SL(2, \mathbb{C})$.
\end{proof}

\begin{theorem}[Cluster decomposition as orbital integrability condition]\label{thm:cluster-integrability}
The cluster decomposition principle plays the same structural role in QFT as the Bianchi identity in GR and the Casimir eigenvalue in QM:

\begin{center}\small
\begin{tabular}{llll}\toprule
\textbf{Theory} & \textbf{Integrability condition} & \textbf{What it forces} & \textbf{Type} \\\midrule
GR & $\nabla_\mu G^{\mu\nu} \equiv 0$ & Uniqueness of field eq. & Identity \\
QM & $C_i = \lambda_i \cdot \mathrm{Id}$ & Uniqueness of wave eq. & Algebraic \\
\textbf{QFT} & \textbf{Cluster decomposition} & \textbf{Existence of local fields} & \textbf{Asymptotic} \\
\bottomrule
\end{tabular}
\end{center}

All three are constraints that make the conservation/unitarity structure self-consistent: Bianchi ensures $\nabla_\mu T^{\mu\nu} = 0$ (energy conservation); Casimir eigenvalues ensure a well-defined Hamiltonian (unitary evolution); cluster decomposition ensures that unitarity is compatible with locality (distant experiments don't interfere).
\end{theorem}


% ============================================================================
\section{The Spin-Statistics Connection}
\label{sec:spinstat}
% ============================================================================

\begin{definition}[Causality condition]\label{def:causality}
For the S-matrix to be both Lorentz-invariant and cluster-decomposing, the fields must satisfy a \textbf{microcausality condition}: for spacelike separation $(x - y)^2 > 0$,
\begin{equation}\label{eq:causality}
[\phi_l(x), \phi_{l'}^\dagger(y)]_\pm = 0,
\end{equation}
where $[\cdot, \cdot]_+$ is the anticommutator and $[\cdot, \cdot]_-$ the commutator. The question is: which sign applies for which spin?
\end{definition}

\begin{theorem}[Spin-statistics --- Weinberg~\cite{weinberg1995}, Pauli~\cite{pauli1940}]\label{thm:spinstat}
For fields transforming in the $(j, 0) \oplus (0, j)$ representation of $\SL(2, \mathbb{C})$:
\begin{enumerate}[label=(\alph*)]
    \item \textbf{Integer spin} ($j \in \{0, 1, 2, \ldots\}$): the causality condition~\eqref{eq:causality} is satisfied only with the \emph{commutator} (minus sign). These fields obey \textbf{Bose--Einstein statistics}.
    \item \textbf{Half-integer spin} ($j \in \{\frac{1}{2}, \frac{3}{2}, \ldots\}$): the causality condition is satisfied only with the \emph{anticommutator} (plus sign). These fields obey \textbf{Fermi--Dirac statistics}.
\end{enumerate}
Using the wrong statistics violates microcausality and hence Lorentz invariance or cluster decomposition.
\end{theorem}

\begin{proof}[Proof outline]
Compute $[\phi(x), \phi^\dagger(y)]_\pm$ using the mode expansion~\eqref{eq:field}:
\begin{equation}
[\phi(x), \phi^\dagger(y)]_\pm = \int \frac{\dd^3p}{(2\pi)^3 2p^0}\left[|u|^2 e^{ip\cdot(x-y)} \pm |v|^2 e^{-ip\cdot(x-y)}\right].
\end{equation}
For spacelike $(x-y)$, one can perform a Lorentz transformation that reverses the sign of $(x-y)^0$. Under this transformation, $e^{ip\cdot(x-y)} \leftrightarrow e^{-ip\cdot(x-y)}$, and the coefficient functions transform as $|u|^2 \to |v|^2 \cdot (-1)^{2j}$ (the factor $(-1)^{2j}$ comes from the representation of the $2\pi$-rotation in $\SL(2, \mathbb{C})$).

For integer $j$: $(-1)^{2j} = +1$, so the commutator ($-$ sign) gives two terms that cancel. The anticommutator ($+$ sign) would give a non-zero result at spacelike separation, violating causality.

For half-integer $j$: $(-1)^{2j} = -1$, so the anticommutator ($+$ sign) gives cancellation. The commutator would violate causality.
\end{proof}

\begin{remark}[Orbital interpretation]
The spin-statistics connection is a \emph{topological constraint on the orbit structure}. The double cover $\SL(2, \mathbb{C}) \to \SO^+(1,3)$ means that half-integer representations pick up a sign under a $2\pi$ rotation. This sign forces anticommutation (to maintain single-valued physics on the orbit space). The statistics is determined by the topology of the representation, not by any dynamical principle.
\end{remark}


% ============================================================================
\begin{remark}[Resolving the Fock-space choice; the Pauli exclusion principle]\label{rmk:fock-resolved}
This settles the symmetric-versus-antisymmetric ambiguity deliberately left open in
Definition~\ref{def:fock}. The $n$-particle space $\mathcal{H}_n$ is the \emph{symmetric} tensor power
$\mathrm{Sym}^n\mathcal{H}_1$ for integer-spin (Bose) fields and the \emph{antisymmetric} power
$\Lambda^n\mathcal{H}_1$ for half-integer-spin (Fermi) fields --- the choice is forced, not postulated, by
microcausality through the $(-1)^{2j}$ factor of Theorem~\ref{thm:spinstat}. The antisymmetry of
$\Lambda^n\mathcal{H}_1$ is precisely the Pauli exclusion principle: a state with two identical fermions in the
same mode vanishes, $a^\dagger(\mathbf{p},\sigma)^2=0$. Thus the exclusion principle --- the backbone of atomic
structure and chemistry --- is a topological consequence of the double cover
$\SL(2,\mathbb{C})\to\SO^+(1,3)$, not an independent postulate.
\end{remark}

% ============================================================================
\section{The CPT Theorem}
\label{sec:cpt}
% ============================================================================

\begin{theorem}[CPT --- L\"uders~\cite{luders1954}, Pauli~\cite{pauli1955}, Jost~\cite{jost1957}]\label{thm:cpt}
Any quantum field theory satisfying Axioms QF1--QF3 (Poincar\'e invariance, unitarity, cluster decomposition) is invariant under the combined operation $\Theta = CPT$:
\begin{equation}
\Theta\,\phi_l(x)\,\Theta^{-1} = \eta_l\,\phi_l^\dagger(-x),
\end{equation}
where $\eta_l$ is a phase and $C$ is charge conjugation, $P$ parity, $T$ time reversal.
\end{theorem}

\begin{proof}[Proof idea]
The Wightman functions (vacuum expectation values of products of fields) are boundary values of analytic functions in the complexified spacetime. The analytic continuation from the forward tube (where $\mathrm{Im}(x_1^0 - x_2^0) > 0$) to the backward tube implements the $PT$ transformation $x \mapsto -x$. Combined with the charge conjugation $C$ (which exchanges the creation and annihilation parts of the field), this gives the CPT transformation~\cite{jost1957,streater1964}.
\end{proof}

\begin{remark}
CPT is a symmetry of the \emph{orbit space}: the complexified Lorentz group $\SO(4, \mathbb{C})$ connects the identity component $\SO^+(1,3)$ to the $PT$-transformed component. CPT invariance is forced by the analytic structure of the representation, not postulated.
\end{remark}


% ============================================================================
\section{Massless Particles and Gauge Invariance}
\label{sec:gauge}
% ============================================================================

\subsection{The little group for massless particles}

For massless particles ($m = 0$), the little group (stabilizer of a reference null momentum $k^\mu = (\omega, 0, 0, \omega)$) is $\ISO(2) = \mathbb{R}^2 \rtimes \SO(2)$, the Euclidean group in 2 dimensions. The physical representations use only the $\SO(2)$ part: discrete helicity $\lambda = \pm s$. The $\mathbb{R}^2$ ``translation'' generators are set to zero to avoid continuous spin representations (which are not observed).

\subsection{Gauge invariance from representation theory}

\begin{theorem}[Gauge invariance for massless spin-1 --- Weinberg~\cite{weinberg1964a}]\label{thm:gauge}
A massless spin-1 field $A_\mu(x)$ has 4 Lorentz components but only 2 physical polarizations (helicity $\pm 1$). Under a Lorentz transformation, $A_\mu$ transforms as:
\begin{equation}\label{eq:gauge-transform}
U(\Lambda)\,A_\mu(x)\,U^{-1}(\Lambda) = \Lambda_\mu{}^\nu A_\nu(\Lambda x) + \partial_\mu\Omega_\Lambda(x),
\end{equation}
where $\Omega_\Lambda$ is an operator-valued function arising from the $\mathbb{R}^2$ part of the $\ISO(2)$ little group. For the S-matrix to be Lorentz-invariant despite this extra term, the coupling of $A_\mu$ to matter must satisfy:
\begin{equation}
\partial_\mu J^\mu = 0 \qquad \text{(current conservation)}.
\end{equation}
This means the interaction is invariant under \textbf{gauge transformations} $A_\mu \to A_\mu + \partial_\mu\Lambda$.
\end{theorem}

\begin{remark}
Gauge invariance is not a postulate but a \emph{consequence} of the representation theory of the little group $\ISO(2)$ for massless spin-1 particles. In the orbital language: gauge transformations are elements of the \emph{kernel} --- transformations that leave the physical observables unchanged. The 2 unphysical polarizations correspond to the 2 gauge degrees of freedom (the $\mathbb{R}^2$ part of $\ISO(2)$).
\end{remark}

\subsection{The loop closes: massless spin-2 implies general relativity}

\begin{theorem}[Massless spin-2 $\Rightarrow$ general covariance --- Weinberg~\cite{weinberg1964b,weinberg1965}]\label{thm:spin2-gr}
For a massless spin-2 field $h_{\mu\nu}(x)$:
\begin{enumerate}[label=(\alph*)]
    \item The Lorentz transformation includes $h_{\mu\nu} \to h_{\mu\nu} + \partial_\mu\xi_\nu + \partial_\nu\xi_\mu$ (linearized diffeomorphism).
    \item Consistency of the S-matrix requires coupling to a conserved, symmetric tensor: $\partial_\mu T^{\mu\nu} = 0$.
    \item Self-coupling (the graviton must couple to its own stress-energy) forces the full nonlinear Einstein equations at the interacting level.
\end{enumerate}
Therefore: any consistent theory of massless spin-2 particles IS general relativity.
\end{theorem}

\begin{proof}[Proof outline]
(a) follows from the $\ISO(2)$ little group analysis, exactly as for spin-1 (Theorem~\ref{thm:gauge}) but with $\mathbb{R}^2$ acting on two indices. (b) follows from requiring the unwanted $\partial_\mu\xi_\nu$ terms to decouple from the S-matrix: this requires $k^\mu \mathcal{M}_{\mu\nu} = 0$ for on-shell amplitudes, which implies conservation $\partial_\mu T^{\mu\nu} = 0$. (c) is Weinberg's argument~\cite{weinberg1964b}: the graviton couples universally to $T^{\mu\nu}$; since the graviton itself carries energy-momentum, it must couple to its own $T^{\mu\nu}$; iterating this self-coupling to all orders produces the full nonlinear Einstein equations $G_{\mu\nu} + \Lambda g_{\mu\nu} = 8\pi G\,T_{\mu\nu}/c^4$.
\end{proof}

\begin{corollary}[Self-consistency loop]\label{cor:loop}
GR is derived by two logically independent paths within the orbital framework:
\begin{itemize}
    \item \textbf{Paper~2:} $\Diff$-equivariance + Noether conservation + second-order $\xrightarrow{\text{Curiel}}$ Einstein.
    \item \textbf{This paper:} Massless spin-2 + unitarity + cluster decomposition $\xrightarrow{\text{Weinberg}}$ Einstein.
\end{itemize}
The two derivations use different axioms and different uniqueness theorems, yet produce the same equation. This is a non-trivial \emph{consistency check} of the orbital framework.
\end{corollary}


% ============================================================================
\section{The Complete Correspondence}
\label{sec:correspondence}
% ============================================================================

\begin{theorem}[Four-theory correspondence]\label{thm:grand-table}
The four instantiations of the orbital pattern are:

\begin{center}\footnotesize
\begin{tabular}{lcccc}\toprule
\textbf{Role} & \textbf{SR (P1)} & \textbf{GR (P2)} & \textbf{QM (P3a)} & \textbf{QFT (P3b)} \\\midrule
Grid cell & $(10, \alpha, 1)$ & $(<\!10, \alpha, \infty)$ & $(10, \beta, 1)$ & $(10, \beta, \infty)$ \\
Group & Poincar\'e & Diff & Bargmann & Poincar\'e on $\mathcal{F}$ \\
Conservation & $\dd Q/\dd t = 0$ & $\nabla_\mu T^{\mu\nu} = 0$ & $U^\dagger U = 1$ & $S^\dagger S = 1$ + cluster \\
Integrability & Mass-shell & Bianchi & Casimir & Cluster decomp. \\
Uniqueness & Max-sym & Curiel & Schur & Weinberg \\
Output & Lorentz + $E^2\!=\!\ldots$ & Einstein & Schr\"odinger & Fields + SS + CPT \\
Lagrangian & Not used & Not used & Not used & Not used \\
\bottomrule
\end{tabular}
\end{center}
\end{theorem}

\begin{remark}[The fifth cell]
The classification grid has a fifth cell: $(\kappa < 10, \beta, \infty)$ = \textbf{quantum gravity}. The orbital framework predicts that this cell must satisfy: (i)~$\Diff$-equivariance (from GR), (ii)~unitarity on curved backgrounds, (iii)~a gravitational cluster decomposition principle, (iv)~classical limit recovering GR, (v)~flat limit recovering QFT, (vi)~massless spin-2 excitations. The missing ingredient is the \emph{uniqueness theorem} for this cell --- the QG analog of Curiel/Schur/Weinberg. Finding it is the central open problem.
\end{remark}


% ============================================================================


% ============================================================================
% PART II: THE UNIFYING OBSERVATION
% ============================================================================

\part*{Part II: The Unifying Observation}

% ============================================================================
\section{The Orbital Framework: Formal Definitions}
\label{sec:framework}
% ============================================================================

\begin{definition}[Physical theory as orbital data]\label{def:theory}
A \textbf{physical theory} in the orbital framework is specified by a quintuple $\mathfrak{T} = (\mathcal{M}, \mathcal{S}, G, \mathcal{C}, \mathcal{U})$: arena $\mathcal{M}$, state space $\mathcal{S}$, symmetry group $G$, conservation principle $\mathcal{C}$, and uniqueness theorem $\mathcal{U}$. The \textbf{equation} $E[\mathfrak{T}]$ is the output of the triple $(G, \mathcal{C}, \mathcal{U})$.
\end{definition}

\begin{definition}[Classification triple]\label{def:triple}
Each theory is characterized by $(\kappa, \rho, \sigma)$: $\kappa \in \{0, \ldots, 10\}$ (kernel dimension), $\rho \in \{\alpha, \beta\}$ (classical/quantum), $\sigma \in \{1, \infty\}$ (single-particle/field).
\end{definition}

\begin{definition}[The classification grid]\label{def:grid}
\begin{equation}\label{eq:grid}
\begin{array}{c|cc}
(\rho, \sigma) \backslash \kappa & \kappa = 10 & \kappa < 10 \\[3pt] \hline
(\alpha, 1) & \textbf{SR} & \text{Classical particle, curved spacetime} \\[3pt]
(\alpha, \infty) & \text{Flat classical field theory} & \textbf{GR} \\[3pt]
(\beta, 1) & \textbf{QM} & \text{QM in curved spacetime} \\[3pt]
(\beta, \infty) & \textbf{QFT} & \textbf{Quantum Gravity} \\[3pt]
\end{array}
\end{equation}
\end{definition}


% ============================================================================
\section{The Unifying Observation}
\label{sec:meta}
% ============================================================================

\begin{observation}[The Equations of Physics from Orbital Structure]\label{thm:meta}
For each of the four occupied cells $(\kappa, \rho, \sigma)$ of the classification grid, the fundamental equation is the unique output of an orbital derivation $(G, \mathcal{C}, \mathcal{U})$:

\begin{center}\footnotesize
\begin{tabular}{cccccc}\toprule
$(\kappa, \rho, \sigma)$ & $G$ & $\mathcal{C}$ & $\mathcal{U}$ & \textbf{Equation} & \textbf{Paper} \\\midrule
$(10, \alpha, 1)$ & $\Poinc(c)$ & $\dd Q/\dd t = 0$ & Max-sym classif. & Lorentz + $E^2\!=\!p^2c^2\!+\!m^2c^4$ & 1 \\[2pt]
$(<\!10, \alpha, \infty)$ & $\Diff(\mathcal{M})$ & $\nabla_\mu\mathcal{E}^{\mu\nu}\!=\!0$ & Curiel/N-S & $G_{\mu\nu}\!+\!\Lambda g_{\mu\nu}\!=\!8\pi T_{\mu\nu}$ & 2 \\[2pt]
$(10, \beta, 1)$ & Bargmann & $U^\dagger U\!=\!1$ & Schur & $i\hbar\partial_t\psi\!=\!H\psi$ & 3 \\[2pt]
$(10, \beta, \infty)$ & $\Poinc$ on $\mathcal{F}$ & $S^\dagger S\!=\!1$ + cluster & Weinberg & Fields + SS + CPT + gauge & here \\[2pt]
$(<\!10, \beta, \infty)$ & $\Diff$ on $\mathcal{F}$ & Unitarity + ??? & ??? & \textbf{Quantum Gravity} & --- \\
\bottomrule
\end{tabular}
\end{center}

No Lagrangian is used as input in any of the four derivations.
\end{observation}

\begin{remark}[Status of the Unifying Observation]
This is an \emph{observation about four rigorously proved instances}, not a theorem about all possible theories. Each row is a theorem (proved in Papers~1--3 and Part~I of this paper). The observation is that all four theorems follow the same three-step pattern. Whether this pattern can be formalized as a single mathematical theorem --- e.g., ``any physical theory satisfying [abstract axioms] must follow this pattern'' --- is an open question. We present it as an empirical observation, not as a proven metatheorem.
\end{remark}


% ============================================================================
\section{The Structural Correspondences}
\label{sec:correspondences}
% ============================================================================

\begin{theorem}[Bianchi--Casimir--Cluster correspondence]\label{thm:correspondence}
The integrability conditions of the four theories are structurally parallel:

\begin{center}\small
\begin{tabular}{lp{2.5cm}p{2.5cm}p{2.5cm}}\toprule
\textbf{Aspect} & \textbf{Bianchi (GR)} & \textbf{Casimir (QM)} & \textbf{Cluster (QFT)} \\\midrule
Object & Tensor $G^{\mu\nu}$ & Operator $C_i$ & S-matrix $S$ \\[3pt]
Condition & $\nabla_\mu G^{\mu\nu} \equiv 0$ & $[C_i, U(g)] = 0$ & $S_{A \cup B} \!\to\! S_A \!\otimes\! S_B$ \\[3pt]
Source & Jacobi identity & Casimir definition & Locality \\[3pt]
Consequence & $\nabla_\mu T^{\mu\nu}\!=\!0$ & $H\!=\!P^2/2m\!+\!E_0$ & Fields $\phi(x)$ exist \\[3pt]
Uniqueness & Curiel: $aG\!+\!bg$ & Schur: $\lambda \!\cdot\! \mathrm{Id}$ & Weinberg: fields forced \\
\bottomrule
\end{tabular}
\end{center}

In each case, the integrability condition is \emph{forced by the group structure}, \emph{reduces the space of possible equations} to finite dimensions, and \emph{guarantees conservation automatically}.
\end{theorem}

\begin{proof}
The structural parallel is verified aspect by aspect.

\textbf{Forced by group structure:} Bianchi follows from the Jacobi identity for covariant derivatives (Paper~2, Theorem~3.1). The Casimir property follows from the definition of the center of the enveloping algebra. Cluster decomposition follows from the locality of the Poincar\'e action on Fock space.

\textbf{Reduces to finite dimensions:} Bianchi + naturality + second-order $\Rightarrow$ $\{aG + bg\}$ (2-parameter family). Casimir on irrep $\Rightarrow$ $\{P^2/2m + E_0\}$ (2-parameter). Cluster + Poincar\'e $\Rightarrow$ fields determined by $(m, s)$.

\textbf{Conservation automatic:} In GR: $\nabla_\mu(aG + bg)^{\mu\nu} = 0$ by Bianchi. In QM: $H = H^\dagger$ guarantees unitarity. In QFT: fields satisfy microcausality by construction.
\end{proof}


% ============================================================================
\section{The Self-Consistency Loop}
\label{sec:loop}
% ============================================================================

\begin{theorem}[Two independent paths to Einstein]\label{thm:loop}
Einstein's equations are derived by two logically independent routes:

\textbf{Path A} (Paper~2): $\Diff$-equivariance + Noether conservation + second-order $\xrightarrow{\text{Bianchi + Curiel}}$ Einstein.

\textbf{Path B} (Part~I, Theorem~\ref{thm:spin2-gr}): Massless spin-2 + unitarity + cluster $\xrightarrow{\text{Weinberg}}$ gauge invariance under linearized diffeomorphisms $\xrightarrow{\text{self-coupling}}$ full nonlinear Einstein.

The two paths use different axioms and different uniqueness theorems, yet produce the same equation.
\end{theorem}

\begin{proof}
Path A: Paper~2~\cite{nembe2026p2}, Theorem~6.1. Path B: Part~I of this paper, Theorem~\ref{thm:spin2-gr}, following Weinberg~\cite{weinberg1964b,weinberg1965} and Deser~\cite{deser1970}.
\end{proof}


% ============================================================================
\section{The Lagrangian as Consequence}
\label{sec:lagrangian}
% ============================================================================

\begin{theorem}[The Lagrangian is always a consequence]\label{thm:lagrangian}
In each theory, a Lagrangian exists \emph{a posteriori}:

\begin{center}\small
\begin{tabular}{llll}\toprule
\textbf{Theory} & \textbf{Equation} & \textbf{Lagrangian} & \textbf{Guaranteeing theorem} \\\midrule
SR & $E^2 = p^2c^2 + m^2c^4$ & $L = -mc^2\sqrt{1 - v^2/c^2}$ & Legendre transform \\
GR & $G + \Lambda g = 8\pi T$ & $\mathcal{L} = R - 2\Lambda$ & Cartan--Vermeil--Weyl \\
QM & $i\hbar\partial_t\psi = H\psi$ & $L = \langle\psi|i\hbar\partial_t - H|\psi\rangle$ & Helmholtz conditions \\
QFT & Fields + interactions & $\mathcal{L}_{\mathrm{QFT}}[\phi]$ & Weinberg reconstruction \\
\bottomrule
\end{tabular}
\end{center}
\end{theorem}

\begin{remark}[The reversal of implication]\label{rmk:reversal-grand}
In every case, the standard approach is: Lagrangian (premise) $\to$ equation $\to$ conservation. The orbital approach reverses this: symmetry + conservation (premise) $\to$ equation $\to$ Lagrangian (consequence). The equations are the same; the logical roles of ``premise'' and ``consequence'' are reversed.
\end{remark}


% ============================================================================
\section{Constraints on Quantum Gravity}
\label{sec:qg}
% ============================================================================

\begin{theorem}[Seven constraints on quantum gravity]\label{thm:qg}
Any theory occupying the cell $(\kappa < 10, \beta, \infty)$ must satisfy:
\begin{enumerate}[label=(QG\arabic*)]
    \item \textbf{Diff-equivariance} (from GR heritage).
    \item \textbf{Unitarity} on curved backgrounds (from QFT heritage).
    \item \textbf{Gravitational cluster decomposition} (multi-particle integrability on curved spacetime).
    \item \textbf{Classical limit:} $\hbar \to 0$ recovers GR.
    \item \textbf{Flat limit:} $\kappa \to 10$ recovers QFT.
    \item \textbf{Graviton content:} massless spin-2 excitations reproducing Newtonian gravity.
    \item \textbf{Uniqueness:} should be determined by a uniqueness theorem (the QG analog of Curiel/Schur/Weinberg). \emph{This is a prediction of the framework, not a derived constraint.}
\end{enumerate}
\end{theorem}

\begin{proof}
(QG1)--(QG2): required for consistency with the $\alpha$ and $\kappa = 10$ limits. (QG3): the multi-particle integrability condition for $\kappa < 10$. (QG4)--(QG5): limiting requirements from neighboring cells. (QG6): from (QG5) + Weinberg's theorem (massless spin-2 $\Rightarrow$ GR). (QG7): a conjecture based on the pattern observed in the other four cells.
\end{proof}

\begin{remark}[Evaluation of existing proposals]
\begin{center}\footnotesize
\begin{tabular}{lccccccc}\toprule
& QG1 & QG2 & QG3 & QG4 & QG5 & QG6 & QG7 \\\midrule
String theory & $\checkmark$ & $\checkmark$ & $?$ & $\checkmark$ & $\checkmark$ & $\checkmark$ & $\times$ (landscape) \\
Loop QG & $\checkmark$ & $?$ & $?$ & $?$ & $?$ & $?$ & $?$ \\
Asymptotic safety & $\checkmark$ & $?$ & $?$ & $\checkmark$ & $?$ & $\checkmark$ & $\checkmark?$ \\
\bottomrule
\end{tabular}
\end{center}
No existing proposal satisfies all seven constraints with certainty.
\end{remark}


% ============================================================================
\section{The Orbital Integrability Conjecture}
\label{sec:oip}
% ============================================================================

\begin{conjecture}[Orbital Integrability Conjecture --- OIC]
\textit{The fundamental equations of physics are the unique integrability conditions that make the orbital structure --- symmetry group, orbits, kernel, and conservation --- self-consistent. The equations are not imposed on nature from outside but emerge from the requirement that symmetries, orbits, and conserved quantities cohere.}
\end{conjecture}

\begin{remark}[Status]
The OIC is supported by four instances (SR, GR, QM, QFT), each rigorously proved. Whether it can be formalized as a single mathematical theorem is an open question. We present it as a conjecture --- a research program, not a proven principle. A possible formalization: in each case, the conserved quantities form a presheaf on the state space, and the equation is the condition that this presheaf extends to a sheaf. We leave this to future work.
\end{remark}

\begin{remark}[Comparison with other principles]
\begin{center}\small
\begin{tabular}{lllc}\toprule
\textbf{Principle} & \textbf{Scope} & \textbf{Status} & \textbf{Lagrangian?} \\\midrule
Least action & Mechanics, field theory & Proven theorem & Yes \\
Equivalence & GR only & Proven & No \\
Gauge principle & Gauge theories & Proven & Yes \\
$\delta Q = TdS$ & GR only & Proven & No \\
\textbf{OIC (this paper)} & \textbf{SR+GR+QM+QFT} & \textbf{Conjecture (4 instances)} & \textbf{No} \\
\bottomrule
\end{tabular}
\end{center}
The OIC has the broadest scope but the weakest epistemic status: it is a conjecture supported by examples, not a proven theorem.
\end{remark}


% ============================================================================
\section{Discussion}
\label{sec:discussion}
% ============================================================================

\subsection{Laws as integrability conditions}

The traditional view holds that the laws of physics are imposed on nature. The OIC suggests an alternative: the laws are internal consistency conditions of the orbital structure. This is an interpretive claim, not a mathematical one; its value lies in its unifying power across the four fundamental theories.

\subsection{The Lagrangian: tool, not foundation}

Theorem~\ref{thm:lagrangian} shows that in every case, the Lagrangian exists as a consequence of the equations. This does not diminish its computational importance --- the Lagrangian generates Feynman rules, conservation laws, and renormalization group flows. It clarifies the logical role: the Lagrangian is a powerful tool, not the foundation.

\subsection{The road to quantum gravity}

The framework predicts that quantum gravity satisfies seven constraints (Theorem~\ref{thm:qg}). The missing piece is the uniqueness theorem for the $(\kappa < 10, \beta, \infty)$ cell. Candidates include: the spectral action (Connes~\cite{connes1996}), the holographic principle, or a yet-unknown classification theorem. Finding this theorem would determine or strongly constrain quantum gravity.


% ============================================================================
\section{Conclusion}
\label{sec:conclusion}
% ============================================================================

We have completed the ``Physics from Orbital Structure'' series by deriving the structure of quantum field theory from the three-step orbital pattern and formalizing the observation that all four fundamental theories follow this pattern.

The key results are: the Bianchi--Casimir--Cluster correspondence (Theorem~\ref{thm:correspondence}); the self-consistency loop via massless spin-2 (Theorem~\ref{thm:loop}); the Lagrangian as consequence (Theorem~\ref{thm:lagrangian}); and seven constraints on quantum gravity (Theorem~\ref{thm:qg}).

We propose the Orbital Integrability Conjecture: the equations of physics are integrability conditions for the self-consistency of orbital structure. This conjecture is supported by four rigorously proved instances but remains open as a general mathematical statement.

\begin{center}
\fbox{\parbox{0.85\textwidth}{\centering\itshape
The equations of physics are not laws imposed on nature\\
but integrability conditions for the self-consistency\\
of its orbital structure.
}}
\end{center}


% ============================================================================
\section*{Acknowledgments}
The author thanks the Roots Insights Laboratory and the National Institute of Management Sciences for support.


% ============================================================================
\begin{thebibliography}{40}

\bibitem{nembe2026p1} J.~Nemb\'e, ``Special relativity from orbital structure,'' Paper~1 of this series (2026).

\bibitem{nembe2026p2} J.~Nemb\'e, ``General relativity from orbital structure,'' Paper~2 of this series (2026).

\bibitem{nembe2026p3a} J.~Nemb\'e, ``Quantum mechanics from orbital structure,'' Paper~3 of this series (2026).

\bibitem{weinberg1995} S.~Weinberg, \emph{The Quantum Theory of Fields}, Vol.~1 (Cambridge Univ.\ Press, 1995).

\bibitem{weinberg1964b} S.~Weinberg, ``Photons and gravitons in S-matrix theory,'' \emph{Phys.\ Rev.}\ \textbf{135}, B1049--B1056 (1964).

\bibitem{weinberg1965} S.~Weinberg, ``Photons and gravitons in perturbation theory,'' \emph{Phys.\ Rev.}\ \textbf{138}, B988--B1002 (1965).

\bibitem{weinberg1964a} S.~Weinberg, ``Feynman rules for any spin.\ II.\ Massless particles,'' \emph{Phys.\ Rev.}\ \textbf{134}, B882--B896 (1964).

\bibitem{wigner1939} E.P.~Wigner, ``On unitary representations of the inhomogeneous Lorentz group,'' \emph{Ann.\ Math.}\ \textbf{40}, 149--204 (1939).

\bibitem{pauli1940} W.~Pauli, ``The connection between spin and statistics,'' \emph{Phys.\ Rev.}\ \textbf{58}, 716--722 (1940).

\bibitem{luders1954} G.~L\"uders, ``On the equivalence of invariance under time reversal and under particle-antiparticle conjugation for relativistic field theories,'' \emph{Kgl.\ Danske Videnskab.\ Selskab, Mat.-fys.\ Medd.}\ \textbf{28}, no.~5, 1--17 (1954).

\bibitem{pauli1955} W.~Pauli, ``Exclusion principle, Lorentz group and reflection of space-time and charge,'' in \emph{Niels Bohr and the Development of Physics}, ed.\ W.~Pauli (Pergamon, London, 1955), pp.~30--51.

\bibitem{jost1957} R.~Jost, ``A remark on the C.T.P.\ theorem,'' \emph{Helv.\ Phys.\ Acta}\ \textbf{30}, 409--416 (1957).

\bibitem{streater1964} R.F.~Streater and A.S.~Wightman, \emph{PCT, Spin and Statistics, and All That} (Benjamin, 1964).

\bibitem{deser1970} S.~Deser, ``Self-interaction and gauge invariance,'' \emph{Gen.\ Rel.\ Grav.}\ \textbf{1}, 9--18 (1970).

\bibitem{wald1986} R.M.~Wald, ``Spin-two fields and general covariance,'' \emph{Phys.\ Rev.\ D}\ \textbf{33}, 3613--3625 (1986).

\bibitem{butcher2009} L.M.~Butcher, M.~Hobson, and A.~Lasenby, ``Bootstrapping gravity: A consistent approach to energy-momentum self-coupling,'' \emph{Phys.\ Rev.\ D}\ \textbf{80}, 084014 (2009).

\bibitem{vermeil1917} H.~Vermeil, \emph{Nachr.\ Ges.\ Wiss.\ G\"ottingen}, 334--344 (1917).

\bibitem{cartan1922} \'E.~Cartan, \emph{Le\c{c}ons sur les invariants int\'egraux} (Hermann, 1922).

\bibitem{navarro2008} J.~Navarro and J.B.~Sancho, ``On the naturalness of Einstein's equation,'' \emph{J.\ Geom.\ Phys.}\ \textbf{58}, 1007--1014 (2008).

\bibitem{curiel2016} E.~Curiel, ``A simple proof of the uniqueness of the Einstein field equation in all dimensions'' (2016).

\bibitem{helmholtz1887} H.~von Helmholtz, \emph{J.\ Reine Angew.\ Math.}\ \textbf{100}, 137--166 (1887).

\bibitem{connes1996} A.~Connes, ``Gravity coupled with matter and the foundation of non-commutative geometry,'' \emph{Comm.\ Math.\ Phys.}\ \textbf{182}, 155--176 (1996).

\bibitem{nembe2026vol17} J.~Nemb\'e, ``Twin General Relativity,'' \emph{ROOTS-INSIGHTS} Vol.~17 (2026).

\end{thebibliography}

\end{document}
